% Generated draft: requires mathematical and source review.
\documentclass[11pt]{article}
\usepackage[margin=25mm]{geometry}
\usepackage{fontspec}
\setmainfont{Latin Modern Roman}
\usepackage{amsmath,amssymb,mathtools,unicode-math}
\setmathfont{Latin Modern Math}
\usepackage{enumitem,needspace,xurl,hyperref}
\setlength{\parindent}{0pt}
\setlength{\parskip}{5pt}
\setlength{\emergencystretch}{3em}
\newcommand{\N}{\mathbb{N}}
\newcommand{\Z}{\mathbb{Z}}
\newcommand{\Q}{\mathbb{Q}}
\newcommand{\R}{\mathbb{R}}
\newcommand{\C}{\mathbb{C}}
\newcommand{\F}{\mathbb{F}}
\newcommand{\A}{\mathbb{A}}
\newcommand{\PP}{\mathbb{P}}
\newcommand{\E}{\mathbb{E}}
\newcommand{\eps}{\varepsilon}
\newcommand{\dd}{\,\mathrm{d}}
\begin{document}
\section*{UnsolvedMath ID: 30003078}
\textbf{Source problem code:} catalogue entry
\subsection{Definitions and mathematical statement}
The supplied question asks which curves containing a finite point set \(Z\) and vanishing to order \(m_Z\) at a general point are unexpected. The record does not define \(m_Z\), specify the ambient space or ground field, or give the allowed degrees. Thus it does not determine a unique mathematical classification problem.

A conditional formulation using the usual dimension comparison for plane curves is as follows. Suppose that the intended setting is \(\PP^2_k\) over an algebraically closed field \(k\), with characteristic still to be specified. Let \(Z\subset\PP^2_k\) be a finite reduced point set, let \(d,m\in\Z_{\geq 0}\), and let \(p\in\PP^2_k\setminus Z\). Write \(\mathcal I_Z\) and \(\mathcal I_p\) for the ideal sheaves of \(Z\) and \(p\), respectively, and define
\[
V_d(Z)=H^0\!\left(\PP^2_k,\mathcal I_Z(d)\right),\qquad
W_{d,m}(Z,p)=H^0\!\left(\PP^2_k,(\mathcal I_Z\cap\mathcal I_p^m)(d)\right).
\]
These are vector spaces of degree-\(d\) forms; a nonzero member of \(W_{d,m}(Z,p)\), modulo nonzero scalar multiplication, defines a curve containing \(Z\) with multiplicity at least \(m\) at \(p\). The expected vector-space dimension after imposing the multiplicity condition is
\[
e_{d,m}(Z)=\max\!\left\{0,\dim_k V_d(Z)-\binom{m+1}{2}\right\}.
\]
In this conditional formulation, the family is unexpected when
\[
\dim_k W_{d,m}(Z,p)>e_{d,m}(Z)
\]
for \(p\) in a nonempty Zariski-open subset of \(\PP^2_k\setminus Z\). Thus unexpectedness refers to an excess dimension of the family, rather than merely to the existence of one curve.

If this plane-curve convention is intended and \(m_Z\) is a specified nonnegative integer associated with \(Z\), the requested classification concerns the degrees \(d\), point sets \(Z\), and corresponding curves for which the displayed inequality holds with \(m=m_Z\). However, neither this interpretation of unexpectedness nor the definition of \(m_Z\) is supplied explicitly. No relation such as \(d=m_Z+1\), no characteristic restriction, and no irreducibility or reducedness condition on the curves can be asserted from the record. These qualifications must be resolved before this conditional formulation can serve as a definitive statement of the source problem.
\subsection{Short English statement}
The question asks when curves through a finite point set, with prescribed multiplicity at a general point, form a larger family than the expected dimension predicts. The displayed plane-curve formulation is conditional: the source record leaves the ambient setting, the meaning of \(m_Z\), the allowed degrees, and the precise classification target unspecified.
\subsection{Sources}
\begin{itemize}
\item Local registry: supplied complete source record, permanent UnsolvedMath ID \(30003078\), title \emph{Classification of Unexpected Curves Through Point Sets}, exact source problem code \texttt{catalogue entry}. The record assigns the year 2016 and supplies no individual proposer. Its statement status is \texttt{corrected\_verified}; the supplied verification describes repairs to notation, typography, and sentence structure without changing mathematical content.
\item Source report: \emph{Mini-Workshop: Arrangements of Subvarieties, and their Applications in Algebraic Geometry} (2016), Oberwolfach Reports, pp. 653--702, DOI \url{https://doi.org/10.4171/owr/2016/14}. The supplied background describes these reports as documenting mathematical workshops and their open-problem discussions. The literature-source list also labels this reference ``OWR report 14/2016.''
\item Supplied literature reference: Cook--Harbourne--Migliore--Nagel, \url{https://doi.org/10.1112/S0010437X18007376}.
\item Supplied literature reference: Di Marca--Malara--Oneto, \url{https://arxiv.org/abs/1804.02730}.
\item Literature assessment dated exactly \texttt{2026-08-21}: the supplied curation labels the problem partially solved and reports a general splitting-type criterion and numerous classified constructions, including supersolvable arrangements and low-degree cases, while reporting no complete geometric classification of all such point sets. This is supplied dated curation derived from the OWR record, not an independently verified literature assessment. The supplied background explicitly requests independent verification of time-sensitive status and cited claims.
\end{itemize}
\end{document}
